\section{Derivative conditions for the kernel}
\label{app:finitepairs}

For each fixed $s>0$, the exponent-pair convention of
\cite[Section 1.1]{TY}, stated in
\cite[Definition~\ref*{six-def:pair}]{Weighted}, specifies a finite
derivative order $d_0$ and a positive tolerance $\epsilon_0$. For a phase
on $I\subset(M,2M]$, the first $d_0$ derivatives must be
relatively $\epsilon_0$-close to
the derivatives of a function with derivative $ym^{-s}$, $y\ge M^s$.
The resulting bound is $(yM^{-s})^kM^\ell$.

The exact reciprocal phase of
\cite[Proposition~\ref*{six-prop:kernel}]{Weighted} has the form $c_Q/m$ with
relative derivative errors $O_j(|\eta|)$ for every fixed $j$.
Here $|c_Q|\asymp FM$. Taking a complex conjugate when required,
use $s=2$ and $y=|c_Q|$. The exponent-pair bound gives
\[
 \sum_{m\in I}e(f(m))\prec F^kM^{\ell-k}
\]
in the range $F\ge c_1M$ for a sufficiently large fixed $c_1$.
Choose $d_0$ first. Since $|\eta|$ is a fixed negative power of $T$,
the derivative errors then satisfy the required bound for large $T$.
The estimate holds on every subinterval. Partial summation gives the
bound for both Sobolev amplitude classes.
The smaller-$F$ region has the separate second-derivative and
small-discrepancy bounds of
\cite[Proposition~\ref*{six-prop:kernel} and
Lemma~\ref*{six-lem:smallF}]{Weighted}.

Bourgain's Theorem 6~\cite{Bou17} gives $(13/84,55/84)$ with
arbitrarily small losses. The finite-order convention and the pair
$A^2(13/84,55/84)=(13/414,359/414)$ are explicit in
\cite[Theorem 1.3]{TY}. All phase parameters are polynomial in
$T$, so the small losses are assigned to the final $\varepsilon$.

For the Heath-Brown pair, write $F=M^\tau$. The desired bound is
\[
 M^{1+\tau/56-31/280+\varepsilon}.
\]
For $\tau\ge31/5$ the length bound suffices. On $1\le\tau\le2$ use
the classical pair $(1/6,2/3)$. On the remaining range use the derivative
estimate of \cite[Theorem 1]{HB17}, rederived from the Vinogradov
mean-value theorem in
\cite[Proposition~\ref*{six-prop:hbderivativeproof}]{Weighted}. With a derivative of order $d$ and
$f^{(d)}\asymp FM^{-d}$, its three normalized costs are
\[
 \frac{\tau-d}{d(d-1)},\qquad
 -\frac1{d(d-1)},\qquad
 -\frac{2\tau}{d^2(d-1)}.
\]
Each is at most $\tau/56-31/280$ on the corresponding interval below:
\[
\begin{array}{c|ccccc}
d&3&4&5&6&7\\\hline
\tau&[2,5/2]&[5/2,17/5]&[17/5,13/3]&[13/3,37/7]&[37/7,31/5].
\end{array}
\]
The comparisons are affine in $\tau$, so their endpoint inequalities
prove the full interval statements. They are included in
\path{checks/finite_derivative_checks.json}.
Thus derivatives through order seven suffice for the entire nontrivial
range of this pair. All derivative constants remain fixed, and both
signs of the Poisson modes satisfy these conditions.
